% =====================================================================
% Utilitários de PLACEHOLDER para o template.
% Pode remover este arquivo (e o \input dele em thesis.tex) quando o
% trabalho estiver com figuras reais.
% =====================================================================

% Caixa cinza tracejada que ocupa o lugar de uma figura ainda não inserida.
% Uso: \figuraplaceholder{largura}{altura}{caminho-sugerido}
%   ex.: \figuraplaceholder{0.6\textwidth}{4cm}{imagens/sua-figura.png}
% Para inserir a imagem real, troque a chamada por:
%   \includegraphics[width=0.6\textwidth]{imagens/sua-figura.png}
\newcommand{\figuraplaceholder}[3]{%
  \begin{tikzpicture}
    \node[draw=gray, fill=gray!12, thick, dashed, rounded corners,
          minimum width=#1, minimum height=#2, align=center,
          text=gray!60, font=\small, inner sep=6pt]
      {[ FIGURA --- substitua por\\ \texttt{\textbackslash includegraphics\{#3\}} ]};
  \end{tikzpicture}%
}
